% =====================================================================
%  Capítulo 5: Discusión y Análisis de Resultados
%  Rúbrica: 15 %. Consistente con el marco teórico; relación resultados –
%  metodología – objetivos; responder TODOS los objetivos.
%  Reestructurado 2026-11-20 según los objetivos aprobados (VF-V2):
%  OE1 caracterización, OE2 diseño depurativo e hidráulico, OE3 estructural.
% =====================================================================
\chapter{Discusión y Análisis de Resultados}
\label{cap:discusion}

\section{Características fisicoquímicas y de producción de las fuentes}
\label{sec:dis-caracterizacion}

Las tres fuentes estudiadas presentan características muy distintas entre sí, lo que justifica tratarlas en líneas independientes.
Las aguas servidas de las oficinas tienen una DQO media de \qty{411}{\milli\gram\per\liter}, dentro del rango informado para las aguas servidas rurales de Chile (\cref{tab:caracteristicas-tipicas})~\parencite{araya2014aguas}, por lo que su carga orgánica es representativa de este tipo de efluente; en cambio, de la casa queda un único dato válido, lo que impide caracterizarla y justifica usar los valores de diseño del proyecto.
Con esos valores, la carga orgánica de la línea doméstica es de \qty{30}{\gram\per\square\meter\per\day} de DQO, y con las concentraciones medidas varía entre unos 17 y \qty{33}{\gram\per\square\meter\per\day} (\cref{tab:validacion-cargas}), por lo que el valor de diseño queda validado.
El amonio medido (1,0 a \qty{1.8}{\milli\gram\per\liter}) queda uno a dos órdenes de magnitud bajo los valores habituales (6 a \qty{60}{\milli\gram\per\liter})~\parencite{araya2014aguas}, y el nitrógeno total no es consistente con el amonio y el nitrato, por lo que el nitrógeno medido no se usa como dato de diseño.
En las oficinas, este nitrógeno bajo se explicaría principalmente por la dilución con el agua del lavamanos y de la cocina; en la casa no tiene una causa confirmada, y como no se midió el caudal de cada artefacto ni la calidad antes y después de la fosa, no es posible afirmarla.

Las aguas de lavado de agroquímicos muestran una variabilidad alta: la DQO de la bodega (\qty{393}{\milli\gram\per\liter}) es cinco veces la de la sala de bombas (\qty{75.7}{\milli\gram\per\liter}), y ambas quedan en el extremo inferior del amplio intervalo reportado para aguas con plaguicidas (150 a \qty{33750}{\milli\gram\per\liter})~\parencite{saleh2020removal}.
Esta variabilidad es propia de un efluente que depende de qué producto se aplica y de cuánta agua se usa en cada lavado, y refuerza el carácter experimental de esa línea.
Además, el volumen de lavado de un predio típico es pequeño, del orden de 5000 a \qty{7000}{\liter} al año~\parencite{braul2018robust}, por lo que la producción real del predio condiciona el caudal que esta línea puede tratar.

El agua del estero Batuco tiene una carga orgánica muy baja y una conductividad eléctrica similar a la de los ríos de la cuenca del Maule~\parencite{dga2004volumen1}; su limitación para el riego es microbiológica, ya que sus coliformes totales la ubican en la clase~4 y su \textit{E.~coli} supera el límite de la clase~1 para coliformes fecales (\cref{tab:clases-calidad}).
Por lo tanto, en esta línea el humedal no se justifica por la remoción de materia orgánica, sino por la reducción de patógenos, para la cual los humedales verticales alcanzan de 2 a 4 unidades logarítmicas~\parencite{langergraber2021vertical}.
Esta eficiencia es teórica y deberá comprobarse durante la operación.

En cuanto a la producción, el caudal de la casa (\qty{0.60}{\cubic\meter\per\day}) se obtuvo de un registro de trece meses con un factor de restitución de \num{0.85}, en el extremo superior del rango de 60 a \qty{85}{\percent} informado en Chile~\parencite{vidal2023desarrollo}, lo que deja el dimensionamiento del lado de la seguridad.
El caudal de las oficinas (\qty{0.15}{\cubic\meter\per\day}) se apoya en solo un mes de registro de usos, por lo que es la estimación con mayor incertidumbre de la línea doméstica.
Para las líneas del estero y de agroquímicos se adoptó la misma carga hidráulica que en la línea doméstica (\qty{750}{\liter\per\day}), porque aplicarles la misma carga orgánica exigiría caudales inviables (unos \qty{28400}{\liter\per\day} en el estero).
Este criterio, coherente con el enfoque sanitario de la estación, hace que cada línea reciba una carga orgánica distinta: \qty{0.8}{\gram\per\square\meter\per\day} en el estero y entre 9,5 y \qty{49}{\gram\per\square\meter\per\day} en agroquímicos según la proporción de AGR1 y AGR2, por lo que en esta última conviene mezclar las fuentes, o completar con agua del estero, para que la DQO de la mezcla no supere unos \qty{240}{\milli\gram\per\liter}.
Como el volumen de lavado de un predio típico es pequeño, el agua del estero cumple además la función de completar el caudal de las otras líneas.

\section{Diseño depurativo y cálculo hidráulico}
\label{sec:dis-hidraulico}

El balance de oxígeno indica que el área activa de \qty{6.00}{\square\meter} por línea aporta \qty{315.0}{\gram\per\day} de O$_2$ frente a una demanda de \qty{111.8}{\gram\per\day} (\cref{tab:balance-oxigeno}), de modo que el lecho puede operar en condiciones aerobias sin aireación forzada~\parencite{platzer1999design,dotro2021humedales}.
Este margen se debe en gran parte al nitrógeno bajo de las concentraciones de diseño: la nitrificación aporta solo \qty{4.7}{\gram\per\day} de la demanda, y con concentraciones de nitrógeno típicas de aguas servidas la demanda aumentaría de forma importante, por lo que el margen debe interpretarse con cautela.

La carga orgánica de \qty{30}{\gram\per\square\meter\per\day} de DQO y la carga hidráulica de \qty{125}{\milli\meter\per\day} superan los máximos de las guías europeas (20 a \qty{27}{\gram\per\square\meter\per\day} y \qty{80}{\liter\per\square\meter\per\day}, \cref{tab:parametros-hcfv}), pero esas guías se desarrollaron para climas con inviernos fríos~\parencite{dotro2021humedales}.
En climas cálidos se han aplicado hasta \qty{200}{\milli\meter\per\day} y \qty{32}{\gram\per\square\meter\per\day} de DBO$_5$~\parencite{platzer2016long}, y diseños recientes en Chile aceptan hasta \qty{35}{\gram\per\square\meter\per\day} de DQO~\parencite{vera2025atrapaniebla}, por lo que las cargas de la Fase~1 quedan dentro de la experiencia regional.
Estas cargas no responden a un criterio conservador, sino a la Fase~1 experimental del proyecto ANID/Anillos/ATE250033, que busca evaluar cuánta agua puede tratarse sin deteriorar la calidad del efluente; el riesgo de colmatación deberá vigilarse en la operación.
Además, los 6 pulsos diarios del balance de oxígeno implican un intervalo de \qty{4}{\hour} entre cargas, dentro del rango de 3 a \qty{6}{\hour} de la guía austríaca, pero menor que las \qty{6}{\hour} que recomienda la guía alemana~\parencite{dotro2021humedales}, por lo que el número de pulsos es un parámetro que convendrá ajustar en la operación.

La configuración en módulos IBC independientes separa físicamente las tres aguas, lo que evita la mezcla entre líneas y permite evaluar el desempeño de cada una por separado, y su instalación sobre el terreno evita excavaciones.
Este tipo de solución en estanques prefabricados es coherente con la experiencia piloto en humedales verticales construidos en recipientes de polietileno~\parencite{ignatowicz2023treatment}, y con la necesidad, señalada para Latinoamérica, de sistemas de distribución presurizada sencillos que faciliten la adopción de los humedales verticales~\parencite{rodriguezdominguez2020constructed}.

En la parte hidráulica, la línea del estero es la más exigente (\cref{tab:bombas-adt}): con PVC de \qty{25}{\milli\meter} en \qty{140}{\meter} y \qty{13.82}{\meter} de desnivel, no alcanza el llenado deseable de \qty{30}{\minute}, pero con la bomba periférica supuesta llena el estanque en unos \qty{36}{\minute}, dentro del límite de \qty{60}{\minute}, sin cambiar el diámetro especificado.
Usar la misma bomba en las tres líneas simplifica la operación y la mantención, porque permite intercambiarla ante una falla.
Por otra parte, el caudal del microrriego (\qty{160}{\liter\per\hour}, \cref{tab:demanda-riego}) es muy inferior al de estas bombas, por lo que requiere otro equipo o un circuito de retorno (\textit{bypass}).
Finalmente, como el trazado de las zanjas se definió en terreno, el diseño entrega los criterios de las zanjas y no su cubicación.

\section{Diseño y cálculo estructural}
\label{sec:dis-estructural}

Frente a las cargas gravitacionales, los perfiles cajón de la plataforma trabajan con razones de demanda a capacidad bajas (alrededor del \qty{9}{\percent} a flexión en la viga principal, \cref{tab:verificacion-estructural}).
Este margen no indica un sobredimensionamiento sin propósito: la condición que controla el diseño es la deformación, porque la plataforma debe mantener nivelados los estanques de almacenamiento y sus conexiones.
En ese sentido, el elemento crítico es el tablero de \qty{15}{\milli\meter}: modelado como un solo vano, no cumplía $L/300$, pero su condición real es la de una lámina continua sobre varias vigas, y con ese modelo cumple con una razón de \num{0.91}.
Como el margen es estrecho y el tablero no cumple la referencia de $L/500$ que la NCh427/1 sugiere para equipos~\parencite{inn2016nch427}, la continuidad de los paños bajo los estanques y el módulo de elasticidad real del tablero fenólico son condiciones del diseño.
Los sifones, en cambio, se ubican en los dosificadores sobre la estructura auxiliar, cuya flecha hace descender el estanque completo sin inclinarlo, aunque sus columnas esbeltas hacen que su estabilidad lateral dependa del arriostramiento y del anclaje.

Frente al sismo, la verificación depende más de cómo se reparte el corte entre los pilares y de la categoría de ocupación que de la versión de la norma (\cref{tab:sismo-zapatas}).
Si los dos pilares de un marco resisten el corte en voladizo, la zapata vuelca con las versiones de 2003 y de 2025; si el tablero y las vigas actúan como diafragma y los seis pilares comparten el corte, con la categoría~I y el suelo~C supuesto el área apoyada es del \qty{91}{\percent}, sobre el mínimo de \qty{80}{\percent} de la NCh2369:2025~\parencite{inn2025nch2369}.
Con la categoría~II el margen desaparece (76 a \qty{79}{\percent}), por lo que la elección de la categoría~I, propia de una estación experimental cuya falla no detiene ninguna operación ni pone en riesgo otras estructuras, es la que permite mantener las zapatas de \qty{1.0}{\meter}.
Por lo tanto, el comportamiento de la plataforma como diafragma es una condición del diseño que debe asegurarse en la conexión del tablero con las vigas.
\revisar{Me falta incorporar la combinación $\num{0.9}D + \num{1.4}E$ de la NCh3171 en el volcamiento~\parencite{inn2017nch3171} y recalcular los pilares con el espectro 2025.}

En el terreno, la presión estática de \qty{24.2}{\kilo\pascal} es muy inferior a la capacidad admisible supuesta de \qty{150}{\kilo\pascal}; según la ecuación de capacidad de carga (\cref{ec:terzaghi-cuadrada}), incluso una arena muy suelta ($\varphi' = \ang{28}$) entregaría unos \qty{86}{\kilo\pascal}~\parencite{das2015fundamentos}.
Por lo tanto, la falta de un estudio de mecánica de suelos no compromete la verificación estática, pero sí la sísmica, en la que las presiones de borde aumentan con la excentricidad.
